\section{Actual native mixed classes and their vanishing graph residue}
\label{sec:136-native-mixed-tor}

\paragraph{Archive and exact scope.}
The uniform conditions are cote136, physical pp.~117 and119,
footers116 and118, scans117 and119; the common-source question is
physical p.~121, footer120, scan121. This calculation concerns two
genuine native degree-four repairs, including their original input
disks. It is not merely the associated graded calculation with the
native differential suppressed. The resulting mixed classes survive
the specified repair in degrees zero and one, but their degree-three
graph obstruction vanishes. No assertion is made about all diagonal
divided-power classes, a later arbitrary horizontal continuation, or
the existence of an integral admissible common source.

Work first over \(R=\mathbb Z_{(2)}\) and
\(S=\Gamma_R(T)\). Put \(M=2^a\), \(a\geq0\), and
\(I=I_M=S_{\geq M}\). Use a graded free resolution
\(P_\bullet\to I\), with augmentation \(\epsilon\) and
differential \(\partial\). Set
\[
                       \mathsf F=S/(2,I).
                                                               \tag{136.NM1}
\]
All gradings and maps below are coefficient-linear.

\paragraph{The native block.}
Begin with the free odd unit disk \(a,b\),
\(da=b\), of forms one and two, both of weight \(m>0\).
Write \(e=ab\), \(z=\gamma_2(b)\); then \(de=2z\).
Adjoin the full degree-four coherent presentation of
ARCH-136-COHERENT-PD-REALIZATION:
\[
\begin{aligned}
dU_0&=\epsilon z,&dU_i&=-U_{i-1}\partial_i,\quad 1\leq i\leq3,\\
dV_0&=2U_0-\epsilon e,&dV_i&=2U_i+V_{i-1}\partial_i,
                                                  \quad1\leq i\leq2,\\
U_3\partial_4&=0,&2U_3+V_2\partial_3&=0.
\end{aligned}                                                   \tag{136.NM2}
\]
Here \(U_i\) has form \(3-i\), \(V_i\) has form \(2-i\),
and the parameter weights are shifted by \(2m\).
Give \(a,b\) auxiliary arity one and every \(U_i,V_i\)
auxiliary arity two. Every displayed equation is homogeneous
for this arity, including its native input term.

Let \(K\) be the arity-two cochain complex of this actual
block. Its form-three term is \(P_0\oplus Se\), its
form-four term is \(Sz\), and the differential between them
is \((u,s)\mapsto\epsilon(u)+2s\). Its terms in forms one
and two are respectively \(P_2\oplus P_1\) and
\(P_1\oplus P_0\). Its term in form zero is
\[
 P_2+\tfrac12Q\subset P_2[1/2],
       \qquad Q=\ker(\partial_2:P_2\to P_1).
                                                               \tag{136.NM3}
\]
This is exactly the bottom-module identification already proved
for (136.NM2), not an assumed free module.

For clarity, the quotient of \(K\) by its native two-term
subcomplex consists just of the coherent parameters. It has
\(H^3=I/2I\) and no other cohomology. Its
bottom differential, in the coordinates (136.NM3), is
\(x\mapsto(2x,\partial_2x)\); its image is exactly the
kernel of the next differential. The higher exactness follows
from exactness of \(P_\bullet\), and its last cokernel is
\(I/2I\). The native two-term block \(Se\xrightarrow{2}Sz\)
is a subcomplex of \(K\). The connecting homomorphism is the
inclusion \(I/2I\to S/2S\), which is injective because
\(S/I\) is free over \(R\). Consequently
\[
                  H^j(K)=0\ (j\ne4),\qquad H^4(K)=\mathsf F.
                                                               \tag{136.NM4}
\]
The common weight shift \(2m\) is understood. Equivalently,
reversing the cochain indices gives a length-four free resolution
of \(\mathsf F\) truncated at its fourth syzygy: terms zero through
three are free, and its last term is that syzygy. Thus (136.NM3)
is also \(\Omega^4\mathsf F\) for this particular resolution.

\paragraph{A truncated tensor calculation.}
If \(Q\) is any projective resolution of a module \(L\)
truncated at homological degree \(r\), with
\(Q_r=\Omega^rL\), then
\[
                 H_j(Q\otimes Q)=\operatorname{Tor}_j(L,L)
                              \quad(0\leq j\leq2r).
                                                               \tag{136.NM5}
\]
Here tensor products are over the coefficient ring.
One proof that also checks the top range is as follows.
Let \(P\to Q\) be the full free resolution and its truncation.
The complex \(Q\otimes P\) computes the derived tensor
product, since \(P\) is a bounded-below complex of free modules.
The map \(Q\otimes P\to Q\otimes Q\) is degreewise onto.
Filter its kernel by the first tensor index \(i\).
For \(i<r\), \(Q_i\) is flat and the corresponding vertical
kernel is acyclic. For \(i=r\), that kernel starts at vertical
degree \(r\); its homology in that degree is zero, since the
image of \(Q_r\otimes P_{r+1}\) is precisely its degree-\(r\)
term. Its only possible homology occurs in vertical degrees
strictly larger than \(r\). The finite first-index filtration
therefore gives zero kernel homology through total degree \(2r\).
The long exact sequence proves (136.NM5). This proof does not
pretend that the last syzygy is flat.

\paragraph{ARCH-136-NATIVE-MIXED-TOR.}
Take the tensor product of two independent native blocks
(136.NM2), and retain their separate auxiliary arities. Its
bidegree \((2,2)\) is exactly \(K\otimes K\). In particular
\[
 H^j((K\otimes K))=
             \operatorname{Tor}_{8-j}^S(\mathsf F,\mathsf F)\langle4m\rangle
                                  \quad(0\leq j\leq8).
                                                               \tag{136.NM6}
\]
There are nonzero order-two classes in degree three in
unbounded weights. These survive the full closed-zeroform
quotient and the subsequent RH1 iteration in the auxiliary
bidegree \((2,2)\), in the homogeneous model over \(S\)
described below. They are not products of extra constants
with positive forms and are not the linear coefficient
syzygies of (136.CG10).

\begin{proof}
The native arity-one piece is the free unit complex
\(Sa\xrightarrow{1}Sb\). Its arity-two piece is precisely
\(K\): the only old monomials are \(e,z\), and a single
coherent parameter already has arity two. There is consequently
no omitted nonlinear monomial in the identification of
bidegree \((2,2)\) with \(K\otimes K\). Applying (136.NM5)
with \(r=4\), and reversing total degree, proves (136.NM6).

We recall the exact input for its nonvanishing.
Put \(t=T^{[M]}\), and let \(C=\Gamma_R(t)\), with
\(t\) of weight \(M\). The coefficient decomposition from
ARCH-136-ACTUAL-TAIL-SYZYGY makes \(S\) finite free and
faithful over \(C\), identifies \(I=SC_+\), and gives
\(\mathsf F=S\otimes_C\mathbb F_2\). Hence
\[
 \operatorname{Tor}_j^S(\mathsf F,\mathsf F)
       =S\otimes_C\operatorname{Tor}_j^C(\mathbb F_2,\mathbb F_2).
                                                               \tag{136.NM7}
\]
The same checked module supplies the free resolution of \(R\)
over \(C\)
\[
 C\otimes\Lambda(e_i:i\geq0)\otimes\Gamma(f_i:i\geq0),
 \quad \partial e_i=x_i,\quad
 \partial f_i=x_i e_i-2\kappa_i e_{i+1},
 \quad x_i=\gamma_{2^i}(t),\quad\kappa_i\in R^\times.
                                                               \tag{136.NM8}
\]
The homological degrees of \(e_i,f_i\) are one and two.
Tensor its complex over \(R\) with the two-term resolution
\(R\xrightarrow{2}R\) of \(\mathbb F_2\), and denote
the new degree-one, weight-zero generator by \(\theta\).
This is a free \(C\)-resolution of \(\mathbb F_2\).
After tensoring with \(\mathbb F_2\) all differentials vanish.
For example
\(\theta\gamma_2(f_i)\) gives a degree-five basis element
of weight \(2^{i+2}M\). Equations (136.NM6)--(136.NM7)
therefore give degree-three classes of weight
\(4m+2^{i+2}M\), for every \(i\geq0\). Their coefficient
summands are copies of \(\mathsf F\), so their order is exactly two.
This uses the already verified resolution (136.NM8), rather
than an extrapolation from a finite computation.

Here is the precise scope of the low-degree repair assertion.
Realize each of the two native disk families on
\(\mathbb Z[\Delta[2]]/\mathbb Z[\operatorname{sk}_1\Delta[2]]\),
and impose (136.NM2) on each relative basis simplex. Faces
and degeneracies carry these basis simplexes to basis simplexes
or to zero, so the native divided-square expressions define
actual simplicial maps. The resulting simplicial algebra \(E\)
has \(E_0=E_1=S\); in degree two there is one block of each
type. It is rationally a free relative disk extension, by the
triangular coordinates already checked for (136.NM2).
It may be mapped to any chosen admissible targets whenever
the indicated coherent images are supplied; the coefficient
augmentation is also always available. Thus RH1 applies with
the target bounds, preserving the two auxiliary arities.

Every componentwise smaller auxiliary bidegree below \((2,2)\)
has zero positive-arity \(H^0\) and \(H^1\): it consists
of \(K\), free arity-one unit complexes, or tensor products
with such an actual unit contraction. In any simplex, multiple
relative basis generators give direct sums of the same complexes
and additional contractible mixed unit pieces. RH1 is therefore
an isomorphism in those smaller bidegrees. In bidegree \((2,2)\)
it adds only zeroform generators and identifying zeroform PD
relations. Any product of such a relation with a positive form
has larger auxiliary bidegree, since the auxiliary-zero part
is just \(S\) in form zero. Thus the modules in positive forms,
and their differentials to higher positive forms, are unchanged
in bidegree \((2,2)\). Only new boundaries of one-forms may
appear. In particular \(H^2\) and \(H^3\) are unchanged
at every step and at the filtered colimit \(E^+\).
The same argument applies to the full closed-zeroform quotient
alone. No lower-arity extra constant exists that could produce
one of the displayed positive-form classes by multiplication.

This assertion is about the specified low-degree repair of
\(E\). A different continuation that introduces new positive
forms is not covered. In particular one cannot automatically
repeat the auxiliary-zero argument after adjoining an arbitrary
old algebra with positive forms in auxiliary arity zero.

The statement also holds integrally, over
\(S_{\mathbb Z}=\Gamma_{\mathbb Z}(T)\). Construct the same
native algebra using any integral free resolution of \(I_M\).
The exactness argument for \(K\) and (136.NM5) are unchanged.
Every group \(\operatorname{Tor}^{S_{\mathbb Z}}_j
(S_{\mathbb Z}/(2,I_M),S_{\mathbb Z}/(2,I_M))\) is annihilated
by two; localization at two is therefore an isomorphism on it.
Localization of a free resolution computes exactly the local
groups above, so their nonvanishing, weights, and annihilation
by \(I_M\) hold integrally. The smaller auxiliary components
still have no \(H^0\) or \(H^1\); this uses the unit
arity-one disks and (136.NM4), not inversion of any odd prime.
The same RH1 construction hence retains the integral mixed
classes while keeping maps to integral targets. The residue
argument below consequently applies over \(\mathbb Z\) as well.
\end{proof}

\paragraph{A coefficient-annihilator test for graph residues.}
For every positive integer \(N\) and every \(r\geq1\),
\[
                 \operatorname{Ann}_{S/NS}(I_M^r)=0.
                                                               \tag{136.NM9}
\]
Here the assertion is also valid over \(\Gamma_{\mathbb Z}(T)\).
For a nonzero element with finite weight support, choose
\(h\geq M\) divisible by \(N d!\), where \(d\) bounds
its weights. Then
\(\binom{h+w}{w}\equiv1\pmod N\) for \(0\leq w\leq d\),
because \(h/i\) is divisible by \(N\) for each \(1\leq i\leq d\).
Multiplication by \(T^{[h]}\) consequently does not kill
the element. Repeating this argument \(r\) times proves
(136.NM9). If \(N\) is a unit, the assertion is vacuous.

Suppose a cochain subcomplex \(B\) of a source in one simplex
restricts to zero on \(G=\operatorname{sk}_1\Delta[n]\),
and has maps into admissible targets whose common graph algebra
is \(F(G)\). Use weights above the fixed vertical comparison
bounds through degree three, the row-extension bounds in forms
zero and one, and the endpoint bound for the common first skeleton.
These are exactly the bounds in ARCH-136-SKELETON-TATE-LIFTS
and are independent of the simplex. In that range there is
a coefficient-linear map
\[
 H^2(P_N(B))\longrightarrow H^1(G;S/NS)                 \tag{136.NM10}
\]
for each target: take an absolute primitive of the image of a
closed Moore-cone pair and restrict that primitive to the graph.
The image pair itself is zero on the graph. The absolute cone
has no degree-one cohomology, so different primitives give the
same graph class. Source boundaries give zero because their
primitives restrict to zero. Addition and coefficient
multiplication commute with this construction, which proves
linearity without choosing a linear family of primitives.
This is the individual relative graph obstruction, not merely
the difference between two targets' obstruction classes.
Since the graph cohomology in (136.NM10) is a finite direct
sum of copies of \(S/NS\), (136.NM9) implies that the map
vanishes on every submodule annihilated by a power of \(I_M\).

\paragraph{ARCH-136-NATIVE-MIXED-RESIDUE.}
Every degree-three Moore pair in auxiliary bidegree \((2,2)\)
of the repaired model \(E^+\), in the uniform target range,
admits primitives in all the chosen targets with zero graph
restrictions. This holds for every modulus \(N\), uniformly
in the simplex. The same conclusion holds after tensoring
\(E^+\) with the free even-disk left Kan envelope of a
common first skeleton and using the multiplication maps into
the targets.

\begin{proof}
Both \(H^2\) and \(H^3\) of the indicated component are
direct sums of the Tor modules in (136.NM6). They are killed
by \(I_M\), as well as by two. The Moore-cone exact sequence
therefore shows that
\[
             I_M^2\,H^2(P_N(E^{+,(2,2)}_n))=0
                                        \quad(N\geq1).
                                                               \tag{136.NM11}
\]
This assertion only uses the two adjacent cohomology groups:
the cone group is an extension of \(H^2/N\) by \(H^3[N]\).
The component restricts to zero on the graph because
\(E^+_0=E^+_1=S\). Equations (136.NM9)--(136.NM11)
therefore kill each target's graph obstruction. A zero graph
class is a cone boundary. Extend a graph primitive of that
boundary to the target simplex, component by component using
the fixed horizontal row bounds, and subtract its cone
differential from the absolute primitive. This produces
primitives with graph values literally zero. These are common
values for all targets and are old source elements, so the
usual graph-relative strict PD attachment is permitted.

For the final assertion let \(L(F)\) be the left Kan envelope
used in ARCH-136-DEGREE-THREE-COMMON-SKELETON. Each \(L(F)_n\)
is a free tensor product of even zero-form unit disks over
\(S\), and has an actual weight-preserving \(S\)-linear
cochain contraction onto \(S\). Thus
\(L(F)_n\otimes E_n^+\) has the same indicated cone
cohomology, without a flatness assumption on \(E_n^+\).
Its first skeleton is exactly \(F\), and the indicated
auxiliary component is still zero there. Multiplication of
the two strict target maps gives its target maps. The preceding
annihilator and graph argument applies unchanged. This is a
factorwise low-repair construction, not a claim about applying
RH1 arbitrarily after tensoring or about later horizontal repairs.
\end{proof}

\paragraph{What remains.}
This calculation both finds genuine nonlinear native classes
and proves that these mixed quadratic classes have zero
degree-three graph obstruction. It does not give an exhaustive
description of the diagonal divided-power components of one
native block. In those components, extending a linear homotopy
as a special derivation gives the arity multiple of a scalar,
which cannot in general be divided by two. No preservation
of a non-coefficient-linear contraction by \(\Gamma\) is
assumed. Further low-cohomology control of the strict attachment
is also needed after the compatible pair is installed.

\paragraph{Verification status.}
The author has read the complete saved proof, including the exact
native arities, the truncated-tensor calculation through its top
degree, the actual minimal-resolution input, every smaller auxiliary
bidegree in the RH1 argument, all-modulus annihilator test, individual
relative graph obstruction, and factorwise common-skeleton scope.
The parent and reviewer1348 each completed a whole saved-proof
read and communicated PASS on 2026-11-22. The parent supplied
the native-arity improvement and the useful factorwise skeleton
extension; these contributions are not recorded as independent
rediscoveries. Reviewer1348 also checked the explicit integral
extension. No formal proof checker or finite computation is asserted.
